% Vietnamese translation for OI Public Library
% Source-PDF: 模板 TEMPLATES/上海交通大学ACM模板.pdf
\documentclass[11pt]{article}
\usepackage{oipl-vn}
\usepackage{fvextra}
\usepackage{longtable}
\usepackage{array}
\usepackage{pdfpages}

\setmonofont{Unifont}[Scale=MatchLowercase]
\DefineVerbatimEnvironment{Code}{Verbatim}{fontsize=\scriptsize,breaklines=true,breakanywhere=true}
\sloppy

\newcommand{\OriginalPDF}{\detokenize{../模板 TEMPLATES/上海交通大学ACM模板.pdf}}
\newcommand{\Topic}[3]{#1 & #2 & #3\\}

\title{Thư viện mã chuẩn}
\author{Hao Yuan\\Khoa Khoa học và Kỹ thuật Máy tính\\Đại học Giao thông Thượng Hải}
\date{23 tháng 10 năm 2003}

\begin{document}
\maketitle
\tableofcontents
\clearpage

\section{Chỉ mục nội dung}

Tài liệu nguồn là một sổ tay mã mẫu dạng ảnh, không có lớp văn bản trích xuất được. Các trang mã trong phần sau được giữ theo đúng bản PDF nguồn; phần dưới đây chuyển các tiêu đề và tên thuật toán trong mục lục sang tiếng Việt để tiện tra cứu.

\subsection{Thuật toán và cấu trúc dữ liệu}

\begin{longtable}{>{\raggedright\arraybackslash}p{0.12\linewidth} >{\raggedright\arraybackslash}p{0.68\linewidth} >{\raggedleft\arraybackslash}p{0.10\linewidth}}
\textbf{Mục} & \textbf{Tên tiếng Việt} & \textbf{Trang}\\
\hline
\Topic{1.1}{Số nguyên độ chính xác cao trong C}{3}
\Topic{1.2}{Số nguyên độ chính xác cao trong C++}{4}
\Topic{1.3}{Số thực độ chính xác cao}{6}
\Topic{1.4}{Lớp phân số}{8}
\Topic{1.5}{Heap nhị phân}{8}
\Topic{1.6}{Cây Winner}{9}
\Topic{1.7}{Cây số}{10}
\Topic{1.8}{Cây đoạn}{10}
\Topic{1.9}{Cây đoạn trong IOI 2001}{11}
\Topic{1.10}{Tập hợp rời nhau}{11}
\Topic{1.11}{Sắp xếp nhanh}{11}
\Topic{1.12}{Sắp xếp trộn}{12}
\Topic{1.13}{Sắp xếp cơ số}{12}
\Topic{1.14}{Chọn phần tử nhỏ thứ k}{12}
\Topic{1.15}{KMP}{12}
\Topic{1.16}{Sắp xếp hậu tố}{13}
\end{longtable}

\subsection{Lý thuyết đồ thị và thuật toán mạng}

\begin{longtable}{>{\raggedright\arraybackslash}p{0.12\linewidth} >{\raggedright\arraybackslash}p{0.68\linewidth} >{\raggedleft\arraybackslash}p{0.10\linewidth}}
\textbf{Mục} & \textbf{Tên tiếng Việt} & \textbf{Trang}\\
\hline
\Topic{2.1}{Đường đi ngắn nhất đơn nguồn: Dijkstra với heap nhị phân}{14}
\Topic{2.2}{Đường đi ngắn nhất đơn nguồn: Bellman--Ford với hàng đợi}{15}
\Topic{2.3}{Cây khung nhỏ nhất: Kruskal}{15}
\Topic{2.4}{Cây khung có hướng nhỏ nhất}{16}
\Topic{2.5}{Ghép cặp cực đại trên đồ thị hai phía}{16}
\Topic{2.6}{Ghép cặp trọng số lớn nhất trên đồ thị hai phía}{17}
\Topic{2.7}{Ghép cặp cực đại trên đồ thị tổng quát}{17}
\Topic{2.8}{Luồng cực đại: Ford--Fulkerson bằng ma trận}{18}
\Topic{2.9}{Luồng cực đại: Ford--Fulkerson bằng danh sách liên kết}{19}
\Topic{2.10}{Luồng cực đại chi phí nhỏ nhất bằng ma trận}{20}
\Topic{2.11}{Luồng cực đại chi phí nhỏ nhất bằng danh sách liên kết}{21}
\Topic{2.12}{Nhận dạng đồ thị dây cung}{21}
\Topic{2.13}{DFS: cầu}{22}
\Topic{2.14}{DFS: điểm cắt}{22}
\Topic{2.15}{DFS: khối song liên thông}{23}
\Topic{2.16}{DFS: sắp xếp tô-pô}{24}
\Topic{2.17}{Thành phần liên thông mạnh}{24}
\end{longtable}

\subsection{Lý thuyết số}

\begin{longtable}{>{\raggedright\arraybackslash}p{0.12\linewidth} >{\raggedright\arraybackslash}p{0.68\linewidth} >{\raggedleft\arraybackslash}p{0.10\linewidth}}
\textbf{Mục} & \textbf{Tên tiếng Việt} & \textbf{Trang}\\
\hline
\Topic{3.1}{Ước chung lớn nhất}{25}
\Topic{3.2}{Định lý số dư Trung Hoa}{25}
\Topic{3.3}{Bộ sinh số nguyên tố}{26}
\Topic{3.4}{Bộ sinh hàm phi Euler}{26}
\Topic{3.5}{Logarit rời rạc}{27}
\Topic{3.6}{Căn bậc hai trong trường hữu hạn}{28}
\end{longtable}

\subsection{Thuật toán đại số}

\begin{longtable}{>{\raggedright\arraybackslash}p{0.12\linewidth} >{\raggedright\arraybackslash}p{0.68\linewidth} >{\raggedleft\arraybackslash}p{0.10\linewidth}}
\textbf{Mục} & \textbf{Tên tiếng Việt} & \textbf{Trang}\\
\hline
\Topic{4.1}{Hệ phương trình tuyến tính trên Z2}{29}
\Topic{4.2}{Hệ phương trình tuyến tính trên Z}{30}
\Topic{4.3}{Hệ phương trình tuyến tính trên Q}{30}
\Topic{4.4}{Hệ phương trình tuyến tính trên R}{31}
\Topic{4.5}{Nghiệm của đa thức}{31}
\Topic{4.6}{Nghiệm của phương trình bậc ba và bậc bốn}{32}
\Topic{4.7}{Biến đổi Fourier nhanh}{33}
\Topic{4.8}{FFT: nhân đa thức}{34}
\Topic{4.9}{FFT: tích chập}{34}
\Topic{4.10}{FFT: đảo bit}{34}
\Topic{4.11}{Quy hoạch tuyến tính: đơn hình nguyên thủy}{36}
\end{longtable}

\subsection{Hình học tính toán}

\begin{longtable}{>{\raggedright\arraybackslash}p{0.12\linewidth} >{\raggedright\arraybackslash}p{0.68\linewidth} >{\raggedleft\arraybackslash}p{0.10\linewidth}}
\textbf{Mục} & \textbf{Tên tiếng Việt} & \textbf{Trang}\\
\hline
\Topic{5.1}{Các thao tác cơ bản}{38}
\Topic{5.2}{Các thao tác mở rộng}{39}
\Topic{5.3}{Bao lồi}{41}
\Topic{5.4}{Đường kính tập điểm}{43}
\Topic{5.5}{Cặp điểm gần nhất}{44}
\Topic{5.6}{Hình tròn}{44}
\Topic{5.7}{Đa giác lồi rỗng lớn nhất}{46}
\Topic{5.8}{Các tâm của tam giác}{47}
\Topic{5.9}{Thể tích đa diện}{48}
\Topic{5.10}{Đường bao của đồ thị phẳng}{49}
\Topic{5.11}{Diện tích hợp các hình chữ nhật}{50}
\Topic{5.12}{Chu vi hợp các hình chữ nhật}{52}
\Topic{5.13}{Đường tròn bao nhỏ nhất}{54}
\Topic{5.14}{Mặt cầu bao nhỏ nhất}{55}
\end{longtable}

\subsection{Bài toán kinh điển}

\begin{longtable}{>{\raggedright\arraybackslash}p{0.12\linewidth} >{\raggedright\arraybackslash}p{0.68\linewidth} >{\raggedleft\arraybackslash}p{0.10\linewidth}}
\textbf{Mục} & \textbf{Tên tiếng Việt} & \textbf{Trang}\\
\hline
\Topic{6.1}{Bộ sinh số Bernoulli}{57}
\Topic{6.2}{Biểu thức Baltic OI 1999}{58}
\Topic{6.3}{Tô màu hạt: lý thuyết Pólya}{58}
\Topic{6.4}{Số Stirling nhị phân}{59}
\Topic{6.5}{Khoảng cách trên mặt hộp}{59}
\Topic{6.6}{Tính biểu thức}{59}
\Topic{6.7}{Cây Descartes}{60}
\Topic{6.8}{Bộ sinh số Catalan}{61}
\Topic{6.9}{Tô màu đa giác đều}{62}
\Topic{6.10}{Đếm cặp nghịch thế}{62}
\Topic{6.11}{Đếm cây}{63}
\Topic{6.12}{Bài toán tám ô}{63}
\Topic{6.13}{Tháp Hà Nội mở rộng}{66}
\Topic{6.14}{Căn bậc hai độ chính xác cao}{66}
\Topic{6.15}{Hình chữ nhật rỗng lớn nhất}{67}
\Topic{6.16}{Chữ số khác 0 cuối cùng của N!}{69}
\Topic{6.17}{Tổ tiên chung thấp nhất}{69}
\Topic{6.18}{Xâu con chung dài nhất}{71}
\Topic{6.19}{Dãy con tăng không giảm dài nhất}{73}
\Topic{6.20}{Nối và tách}{73}
\Topic{6.21}{Hình vuông ma thuật}{74}
\Topic{6.22}{Cây tìm kiếm nhị phân tối ưu}{75}
\Topic{6.23}{Xếp các hình chữ nhật: cắt hình chữ nhật}{75}
\Topic{6.24}{Xếp các hình chữ nhật: O(N\textasciicircum{}2)}{76}
\Topic{6.25}{Parliament}{76}
\Topic{6.26}{Bộ sinh pi}{76}
\Topic{6.27}{Cây Plant: lặp}{77}
\Topic{6.28}{Cây Plant: cây đoạn}{77}
\Topic{6.29}{Truy vấn cực đại trên đoạn}{78}
\Topic{6.30}{Bài toán người bán hàng}{79}
\Topic{6.31}{Chiều cao cây}{80}
\Topic{6.32}{Biểu diễn chu trình nhỏ nhất}{81}
\Topic{6.33}{Clique cực đại}{82}
\Topic{6.34}{Ma trận con tối đại không cấm}{83}
\Topic{6.35}{Bài toán hai dây chuyền cực đại}{83}
\Topic{6.36}{Bài toán N quân hậu}{85}
\Topic{6.37}{Bộ sinh dãy De Bruijn}{85}
\Topic{6.38}{ZOJ 1482: phân hoạch}{86}
\end{longtable}

\clearpage
\section{Các trang mã nguồn}

\noindent Phần này giữ nguyên các trang của sổ tay mã chuẩn để bảo toàn từng khối mã, định danh và chú thích trong mã.

\includepdf[pages=-,pagecommand={},fitpaper=false]{\OriginalPDF}

\end{document}
